% ---
% filename : nibt_installations_20260517_900_sectionnement_abonne.tex
% id : nibt_installations_20260517_900
% title : "Sectionnement et disjoncteurs d'abonné"
% domain : "Nibt"
% subdomain : "Canalisations"
% tags : ["nibt", "sectionnement", "disjoncteur", "sécurité", "dt"]
% difficulty : 2
% time_solve : 15
% has_octave : false
% author : "om"
% ---
\documentclass[a4paper,11pt,answers,addpoints]{exam}
% ---------------------------------------------------------
% LANGUE, ENCODAGE, FONTS
% ---------------------------------------------------------
\usepackage[utf8]{inputenc}
% ---------------------------------------------------------
% GÉOMÉTRIE ET MISE EN PAGE
% ---------------------------------------------------------
\usepackage[left=1.7cm,right=1.5cm,top=2cm,bottom=1cm]{geometry}
\usepackage{microtype}
\usepackage{float}
\usepackage{needspace}

% ---------------------------------------------------------
% MATHÉMATIQUES
% ---------------------------------------------------------
\usepackage{amsmath, amssymb, bm, esvect}

\newcommand{\V}{\overrightarrow}
\def\vect#1{\overrightarrow{\strut#1}}
\providecommand{\Degrees}[0]{\ensuremath{^\circ}}

% ---------------------------------------------------------
% UNITÉS ET GRANDEURS (SIunitx)
% ---------------------------------------------------------
\usepackage{siunitx}
\sisetup{output-decimal-marker={,}}
\DeclareSIUnit \var {var}
\DeclareSIUnit \VA {VA}
\DeclareSIUnit \kVA {kVA}
\DeclareSIUnit[quantity-product={}] \degree{\text{\textdegree}}

% Permet la césure dans les nombres avec unités (évite les Underfull \hbox)
%\AllowBreakAtQQ

% ---------------------------------------------------------
% GRAPHIQUES : TikZ + CircuitikZ (sans tkz-fct qui cause des erreurs spath3)
% ---------------------------------------------------------
\usepackage{tikz}
\usetikzlibrary{
	arrows.meta,
	intersections,
	decorations.pathreplacing,
	positioning
}

\usepackage[european,straightvoltages,EFvoltages]{circuitikz}
\usepackage{pgfplots}
\pgfplotsset{compat=1.12}

% Symbole étoile-triangle personnalisé
\newcommand{\wye}{%
	\mathbin{\tikz[x=1ex,y=1ex]{%
			\draw[line width=.1ex] (0,0)--(45:1)--++(-45:1) (45:1)--++(0,1);
	}}%
}

% ---------------------------------------------------------
% TABLEAUX
% ---------------------------------------------------------
\usepackage{tabularx}
\usepackage{tabu}

% ---------------------------------------------------------
% HYPERLIENS
% ---------------------------------------------------------
\usepackage[colorlinks=true,
menucolor=red,
breaklinks=true,
urlcolor=blue,
linkcolor=blue]{hyperref}

% ---------------------------------------------------------
% COMMANDES PERSONNELLES
% ---------------------------------------------------------
\newcommand\nomauteur{bdminasmorgul@protonmail.com}
\newcommand\dateTe{18.05.2026}
\newcommand\brancheTe{DT}

% ---------------------------------------------------------
% STYLE DE L'EXERCICE
% ---------------------------------------------------------
\pagestyle{headandfoot}
\firstpageheadrule
\runningheadrule

% En-tête avec largeur fixe pour éviter les dépassements
\firstpageheader{\brancheTe\ (\nomauteur)}{Élève : }{\makebox[3.5cm][r]{\dateTe\ \thepage/\numpages}}
\runningheader{\brancheTe\ (\nomauteur)}{Élève : }{\makebox[3.5cm][r]{\dateTe\ \thepage/\numpages}}

% Pied de page
\newcommand{\totalpointtts}{%
	\begin{tikzpicture}[remember picture, overlay]
		\node[anchor=south west, inner sep=0pt, outer sep=0pt]
		at ([xshift=0.5cm,yshift=0.5cm]current page.south west)
		{\rotatebox{90}{Total des points : \numpoints}};
	\end{tikzpicture}%
}

\firstpagefooter{\hspace*{\fill}\totalpointtts}{}{}
\runningfooter{\hspace*{\fill}\totalpointtts}{}{}

% Réglages exam
\printanswers
\addpoints
\pointsinmargin
\bracketedpoints

% ---------------------------------------------------------
% LISTINGS (code)
% ---------------------------------------------------------
\usepackage{xcolor}
\usepackage{listings}

\lstdefinestyle{octavestyle}{
	language=Matlab,
	basicstyle=\ttfamily\small,
	keywordstyle=\color{blue},
	commentstyle=\color{green!50!black},
	stringstyle=\color{red},
	stepnumber=1,
	frame=single,
	breaklines=true,
	breakatwhitespace=true,
	tabsize=4
}

\usepackage{enumitem}
\setlist[enumerate,1]{label=\alph*), ref=\alph*)}
\newenvironment{explication}{%
	\par\noindent\textbf{Explication. }\itshape
}{\par\normalfont}

\begin{document}
	\unframedsolutions
	\pointsinmargin
	\bracketedpoints
	
\section*{Préparation DT PQ CFC ELMO,IELE,PELE}
\begin{questions}
\question
En tant qu'installateur-électricien, vous réalisez le tableau de distribution pour les nouveaux appartements d'un immeuble locatif.

\begin{parts}
	\part[3] Quelle est la particularité mécanique que doivent présenter les pôles des conducteurs de phase d'un disjoncteur utilisé comme coupe-surintensité d'abonné~?
	\part[3] Dans quel cas précis est-il autorisé d'utiliser des disjoncteurs de canalisation à 2 pôles (L + N) pour la protection d'un abonné sans contrevenir à la règle de non-couplage mécanique~?
	\part[3] À quel seul et unique endroit est-il autorisé d'installer un sectionneur dans le conducteur PEN~?
	\part[3] En dehors du coupe-surintensité général et de celui de l'abonné, citez un autre emplacement où une installation doit obligatoirement pouvoir être sectionnée~?
\end{parts}

\begin{solution}
	Le sectionnement et la protection des circuits d'abonnés répondent à des règles strictes pour garantir l'indépendance des phases et la sécurité des interventions :
	
	\begin{itemize}
		\item \textbf{Couplage mécanique} : Les pôles des conducteurs de phase ne doivent pas être couplés mécaniquement. Cela permet d'éviter qu'un défaut sur une phase n'entraîne la coupure immédiate des autres phases de l'abonné si cela n'est pas souhaité (NIBT 4.6.2.2.2).
		\item \textbf{Exception L+N} : Il est autorisé d'utiliser des disjoncteurs L+N si les conducteurs neutres respectifs sont montés en parallèle côté entrée et sortie, permettant une utilisation par des personnes ordinaires sans outil (NIBT 4.6.2.2.2).
		\item \textbf{Sectionnement du PEN} : Le sectionneur de neutre ne peut être installé dans le conducteur PEN qu'au niveau du coupe-surintensité général (HAK : Hausanschlusskasten) (NIBT 4.6.1.2.1.1).
		\item \textbf{Points de sectionnement} : Un sectionnement doit être prévu sur tous les circuits partant d'un ensemble d'appareillage (NIBT 4.6.2.1.3).
	\end{itemize}
\end{solution}
\end{questions}
\end{document}
